\cotefile{78}
\scanpage{89}
On retrouve encore une fois que \deleted{\unsure{Cela|le}} la donnée ici coïncide avec ce que nous avions appelé (dans le cas des \unsure{``polygônes|polygônes|n-polygônes} réguliers avec $n$ premier à la car. $p$) ``l'invariant numérique'' dudit. Cependant ici le cas \unsure{d'|de'}une racine double
\[
\unsure{\zeta|\rho}=\zeta' \quad \text{i.e.}\quad \zeta^{2}=1 \quad \text{i.e.}\quad \zeta=\pm1 \tag*{(35)}
\]
n'est pas exclu, il correspond au cas
\[
\alpha=\pm2 \qquad \text{ou encore}\quad \Delta(\alpha)=0 \tag*{(36)}
\]
($\alpha=2$ correspond à $\zeta=1$, $\alpha=-2$ à $\zeta=-1$, les deux cas venant se confondre \unsure{sss|ssi} $p=2$).

\unsure{Celà|Cela} tient compte précisément des 2 seules valeurs de $\zeta$ que ns avions été obligés d'exclure, pour définir nos repr. géom. can. $\varphi^\zeta=(\varphi_{\unsure{s|5|0}}^\zeta,\varphi_{\unsure{a|0|o|1}}^\zeta)$ de $\Pi_\infty$ ! [Le cas $\zeta=1$, $\alpha=2$ est également distingué \deleted{\unsure{en}} comme le seul cas où il n'y ait pas de centre à distance finie .... ]
% archive sheet number: 88
